% Unidad R565: introducción, p. 5, dedicatoria y requisitos a)--d).
% Continuación: intro-fr.tex, párrafo «Ce n'est pas le lieu...».
\section*{Introducción}

\begin{center}
\emph{A Oscar Zariski y André Weil.}
\end{center}

\oldpage[I]{5}

Este trabajo, y los numerosos trabajos que han de seguirlo, están destinados a formar un tratado sobre los fundamentos de la geometría algebraica.
En principio no presuponen ningún conocimiento particular de esta disciplina; incluso se ha comprobado que tal conocimiento, pese a sus evidentes ventajas, podía a veces resultar perjudicial para quien desea familiarizarse con el punto de vista y las técnicas expuestos aquí, por el hábito de adoptar de manera demasiado exclusiva el punto de vista birracional que dicho conocimiento lleva consigo.
En cambio, supondremos que el lector conoce bien los temas siguientes:

\begin{enumerate}
  \item[a)] El \emph{álgebra conmutativa}, tal como se expone, por ejemplo, en los volúmenes de los \emph{Éléments} de N.~Bourbaki que están en preparación (y, hasta que aparezcan esos volúmenes, en Samuel-Zariski~\cite{I-13} y Samuel~\cite{I-11},~\cite{I-12}).
  \item[b)] El \emph{álgebra homológica}, para la cual remitimos a Cartan-Eilenberg~\cite{I-2} (citado como~(M)) y Godement~\cite{I-4} (citado como~(G)), así como al reciente artículo de A.~Grothendieck~\cite{I-6} (citado como~(T)).
  \item[c)] La \emph{teoría de haces}, cuyas referencias principales serán~(G) y~(T); esta teoría proporciona el lenguaje indispensable para interpretar en términos «geométricos» las nociones esenciales del álgebra conmutativa y para «globalizarlas».
  \item[d)] Por último, será útil que el lector esté familiarizado en cierta medida con el \emph{lenguaje funtorial}, que se empleará constantemente en este tratado y para el cual podrá consultar~(M),~(G) y, sobre todo,~(T); los principios de este lenguaje y los principales resultados de la teoría general de los funtores se expondrán con mayor detalle en una obra que los autores de este tratado están preparando.
\end{enumerate}

\bigskip
\begin{center}*\quad*\quad*\end{center}
\bigskip
